\section{Derivations and consistency checks}\label{sec:derivations}

\subsection{Optical path, branch selection and Cartesian curvature}
For an arbitrary virtual tangential displacement $\bm\tau_t$ at a surface
point, stationarity of the signed optical path reads
$\delta\Psi=(n_o\bm d_o-n_i\bm d_i)\cdot\bm\tau_t=0$. This proves
\cref{eq:fermat-snell}. Conversely, on a connected regular patch, tangential
Snell refraction into the prescribed point implies $\Gs\Psi=0$ and hence a
constant signed path. The equality of tangential components does not select
the direction of the transmitted normal component, which is why the branch
tests accompanying \cref{eq:snell-vector} are indispensable.

For a nonzero level gradient $\bm g_\Psi$, differentiating its unit direction
gives
\begin{equation}
 \nabla\cdot\frac{\bm g_\Psi}{|\bm g_\Psi|}
 =\frac{\tr\bm H_\Psi}{|\bm g_\Psi|}
  -\frac{\bm g_\Psi^T\bm H_\Psi\bm g_\Psi}{|\bm g_\Psi|^3}.
 \label{eq:implicit-curvature-proof}
\end{equation}
The Hessian of a distance from a point is
$(\Id-\bm u\otimes\bm u)/R$, where $R$ is that distance and $\bm u$ its radial
unit vector. These two identities establish
\cref{eq:cartesian-traction-explicit} without a paraxial approximation.

The rational Cartesian parameterization follows from solving
\cref{eq:gots-polynomial} as a quadratic in $z$. Its discriminant is
\begin{equation}
 (1+\mathsf S\varrho^2)^2
 -\mathsf O\mathsf G(\mathsf O+\mathsf T\varrho^2)\varrho^2
 =1+(2\mathsf S-\mathsf O^2\mathsf G)\varrho^2,
 \label{eq:gots-discriminant}
\end{equation}
because $\mathsf S^2=\mathsf O\mathsf G\mathsf T$. The root vanishing at the
vertex is rationalized to obtain \cref{eq:gots-parametric}. This cancellation
does not excuse skipping the original unsquared path equation.

\subsection{Area, curvature and the pressure multiplier}
Define a dimensionless perturbation parameter $\varepsilon$ and a displaced
surface by $\bm x_\varepsilon=\bm x+\varepsilon\eta\nn$. A dot in this
subsection denotes differentiation with respect to $\varepsilon$ at zero.
The elementary geometric variations are
\begin{equation}
 \dot{\nn}=-\Gs\eta,\qquad
 \dot{(\dd A)}=\kappa\eta\dd A,\qquad
 \dot\kappa=-\Delta_\Gamma\eta-|\bm B_\Gamma|^2\eta.
 \label{eq:geometric-variations}
\end{equation}
For a sphere, $\kappa=2/R_s$ and a uniform radial displacement gives
$\dot\kappa=-2\eta/R_s^2$, in agreement with direct differentiation. Domain
transport gives $\delta\int_{\Omega^-}\Phi\dd V=\int_\Gamma\Phi\eta\dd A$.
Using the functional $E-\lambda_V(V-V_0)$ with acoustic virtual work then
yields \cref{eq:required-traction}. At equilibrium, varying its residual gives
\cref{eq:effective-stiffness}, including $-\delta\lambda_V$ before projection.
Derivative terms from the surface measure multiply the vanishing residual;
they cannot be discarded indiscriminately away from equilibrium.

\subsection{Dissipation after actuation has ended}
Consider the restricted incompressible, constant-property, unforced problem
with fixed outer walls, no incoming flux, continuous interfacial velocity and
either a pinned contact line or a consistent dissipative wetting/slip law.
Define total kinetic energy
$K_f=\sum_{\ell\in\{-,+\}}\int_{\Omega^\ell}
\rho_\ell|\bm U_\ell|^2/2\dd V$. Let $\mathcal D_{\rm wall}\geq0$ denote
wall-slip and contact-line dissipation; for the line law
\cref{eq:contact-friction}, its line contribution is
$\int_{\mathcal C}\zeta_{\mathcal C}V_{\mathcal C}^2\dd s$. Multiplying
momentum by velocity, integrating by parts and using the interface balance
gives
\begin{equation}
 \frac{\dd}{\dd t}(K_f+E)
 =-\sum_\ell\int_{\Omega^\ell}2\mu_\ell
                  \bm D(\bm U_\ell):\bm D(\bm U_\ell)\dd V
   -\mathcal D_{\rm wall}\leq0.
 \label{eq:passive-energy-decay}
\end{equation}
The surface work cancels the rate of capillary and wetting energy, and gravity
cancels the rate of gravitational energy. The identity provides a check on a
future flow solver and explains the loss of pulse energy. It does not prove
global convergence, exclude metastable equilibria, or apply unchanged during
heating, mass exchange, curing or ongoing wave forcing.

\subsection{Quadratic sensitivities and local load authority}
At fixed geometry define the real vector
$\bm u=(\Rea\bm w,\Ima\bm w)^T\in\R^{2N}$. For a Hermitian $\bm K_i$,
variation of the generalized force is
\begin{equation}
 \delta F_i=2\Rea[(\bm K_i\bm w)^\dagger\delta\bm w],\qquad
 (\mathsf J_F)_{i,:}
    =2\big(\Rea(\bm K_i\bm w)^T,\Ima(\bm K_i\bm w)^T\big),
 \label{eq:quadratic-jacobian}
\end{equation}
where $\mathsf J_F$ is the real load Jacobian. For a nonzero drive, a common
phase shift has tangent $\delta\bm w=\mathrm i\bm w\,\delta\varphi$ with
real phase increment $\delta\varphi$, and leaves every $F_i$ unchanged. Thus
\begin{equation}
 \rank\mathsf J_F\leq\min(M,2N-1).
 \label{eq:local-rank-bound}
\end{equation}
Consequently a finite coherent array cannot independently prescribe an
unlimited number of infinitesimal load components at a fixed geometry. At zero
drive the first derivative vanishes because the load is quadratic; finite
actuation can still deform the surface. Bounds further restrict the available
directional changes.

For a family of targets, changes in $\Gamma$ also change $\bm K_i$ and $b_i$.
Linearizing the target-load equations gives
\begin{equation}
 (\mathsf J_F\delta\bm u)_i
       =\delta b_i-\bm w^\dagger(\delta\bm K_i)\bm w.
 \label{eq:target-to-drive-tangent}
\end{equation}
The right side includes the complete curvature, wave and surface-measure
derivatives. It must lie in the admissible Jacobian range for that local
continuation. Small singular values after physical scaling imply sensitive
control demands. This is a local test of load authority, distinct from
dynamic controllability and from a global existence result. The four-parameter
Cartesian family is restricted enough that the dimensional bound alone need
not obstruct it; range compatibility, limits and stability still matter.

\Needspace{24\baselineskip}
\subsection{Normalization and sign audit}
The following dimensions provide checks for subsequent implementations.
\begin{center}
\begin{tabular}{@{}ll@{}}
\toprule Quantity & SI units \\\midrule
$\Psi,W,R_o,R_i,\eta$ & $\mathrm m$\\
$\kappa,\bm H_\Psi$ & $\mathrm{m^{-1}}$\\
$\sigma\kappa,\Delta\rho\Phi,\lambda_V,\Pi,\bm R$ & $\mathrm{Pa}$\\
$\bm w$; $\bm W$ & $\mathrm{m\,s^{-1}}$; $\mathrm{m^2\,s^{-2}}$\\
$\bm K$; $\bm K_i$ & $\mathrm{kg\,m^{-3}}$; $\mathrm{kg\,m^{-1}}$\\
$F_i,b_i$; $\bm R_P$ & $\mathrm N$; $\mathrm{kg\,s^{-1}}$\\
$E,K_f$; $P_{\rm in}$ & $\mathrm J$; $\mathrm W$\\
$J_W,J_d,J_s,J_{\rm spot},J$; $\ell$ & dimensionless; $\mathrm{s^{-1}}$\\
\bottomrule
\end{tabular}
\end{center}
As a radiation-flux check, a locally progressive plane wave with propagation
unit vector $\bm d$ obeys $\bm V=P\bm d/(\rho c)$. \Cref{eq:radiation-flux} then gives
$\bm R=|P|^2\bm d\otimes\bm d/(2\rho c^2)$: the momentum flux has the sign
of propagation. For the lens normal convention, the hydrostatic balance is
$p^- -p^++\Pi=\sigma\kappa$, consistent with
\cref{eq:required-traction}. Neither check uses a numerical solution or
apparatus-specific values.
